% Part: first-order-logic
% Chapter: natural-deduction
% Section: quantifier-rules

\documentclass[../../../include/open-logic-section]{subfiles}

\begin{document}

\olfileid{fol}{ntd}{qrl}

\olsection{量词规则}

\subsection{$\lforall$ 的规则}

\begin{defish}
\AxiomC{$!A(a)$}
\RightLabel{\Intro{\lforall}}
\UnaryInfC{$\lforall[x][\Atom{!A}{x}]$}
\DisplayProof
\hfill
\AxiomC{$\lforall[x][\Atom{!A}{x}]$}
\RightLabel{\Elim{\lforall}}
\UnaryInfC{$!A(t)$}
\DisplayProof
\end{defish}

在 $\lforall$ 的规则中，$t$ 是闭项（不含任何变元的项），而 $a$ 是常项符号，它不出现在结论~$\lforall[x][!A(x)]$ 中，也不出现在以 $!A(a)$ 为前提的推导中任何未解除的假设里。我们称 $a$ 为 \Intro{\lforall} 推理的\emph{特征变元}。\footnote{尽管上面规则中的 $a$ 是常项，我们仍使用“特征变元”这个术语。这是历史原因。}

\subsection{$\lexists$ 的规则}

\begin{defish}
\AxiomC{$\Atom{!A}{t}$}
\RightLabel{\Intro{\lexists}}
\UnaryInfC{$\lexists[x][\Atom{!A}{x}]$}
\DisplayProof
\hfill
\AxiomC{$\lexists[x][\Atom{!A}{x}]$}
\AxiomC{[$\Atom{!A}{a}$]$^n$}
\DeduceC{$!C$}
\DischargeRule{\Elim{\lexists}}{n}
\BinaryInfC{$!C$}
\DisplayProof
\end{defish}

同样，$t$ 是闭项，$a$ 是常项符号，它不出现在前提~$\lexists[x][!A(x)]$ 或结论~$!C$ 中，也不出现在以这两个前提结尾的推导中任何未解除的假设里（假设 $!A(a)$ 本身除外）。我们称 $a$ 为 \Elim{\lexists} 推理的\emph{特征变元}。

特征变元既不出现在前提中，也不出现在通往 \Intro{\lforall} 或 \Elim{\lexists} 推理之前提的推导中任何未解除的假设里，这一条件被称为\emph{特征变元条件}。

\begin{explain}
回想一下我们的约定：当 $!A$ 是变元~$x$ 在其中自由的公式时，我们将其记为~$!A(x)$。在同一语境下，$!A(t)$ 则是~$\Subst{!A}{t}{x}$ 的简写。因此，我们也可以将 $\Intro\lexists$ 规则写为：
\begin{prooftree}
\AxiomC{$\Subst{!A}{t}{x}$}
\RightLabel{\Intro{\lexists}}
\UnaryInfC{$\lexists[x][!A]$}
\end{prooftree}
注意，$t$ 可能已经在~$!A$ 中出现，例如 $!A$ 可能是~$\Atom{\Obj P}{t,x}$。因此，从~$\Atom{\Obj P}{t,t}$ 推出 $\lexists[x][\Atom{\Obj
P}{t,x}]$ 是~$\Intro\lexists$ 的一个正确应用——你可以用约束变元~$x$“替换”前提中~$t$ 的一次或多次（且不必是全部）出现。然而，$\Intro\lforall$ 和~$\Elim\lexists$ 中的特征变元条件要求常项符号~$a$ 不出现在~$!A$ 中。所以，你不能使用~$\Intro\lforall$ 从 $\Atom{\Obj P}{a,a}$ 正确地推出 $\lforall[x][\Atom{\Obj P}{a,x}]$。
\end{explain}

\begin{explain}
在 \Intro{\lexists} 和 \Elim{\lforall} 中没有任何限制，项~$t$ 可以是任意项，因此我们无需考虑任何条件。另一方面，在 \Elim{\lexists} 和 \Intro{\lforall} 规则中，特征变元条件要求常项符号~$a$ 不出现在结论或任何未解除的假设中。该条件对于确保系统的可靠性是必要的，即系统只能从未解除的假设中推导出作为其逻辑后承的语句。如果没有这个条件，以下推导就会被允许：
\begin{prooftree}
  \AxiomC{$\lexists[x][!A(x)]$}
  \AxiomC{$\Discharge{!A(a)}{1}$}
  \RightLabel{*\Intro{\lforall}}
  \UnaryInfC{$\lforall[x][!A(x)]$}
  \RightLabel{\Elim{\lexists}}
  \BinaryInfC{$\lforall[x][!A(x)]$}
\end{prooftree}
然而，$\lexists[x][!A(x)] \Entails/ \lforall[x][!A(x)]$。

由于量词的消去规则只允许将闭项代入变元，因此，任何能从一个语句集合中推导出的公式本身也是语句。
\end{explain}

\end{document}